%Version 3.1 December 2024
% See section 11 of the User Manual for version history
%
%%%%%%%%%%%%%%%%%%%%%%%%%%%%%%%%%%%%%%%%%%%%%%%%%%%%%%%%%%%%%%%%%%%%%%
%%                                                                 %%
%% Please do not use \input{...} to include other tex files.       %%
%% Submit your LaTeX manuscript as one .tex document.              %%
%%                                                                 %%
%% All additional figures and files should be attached             %%
%% separately and not embedded in the \TeX\ document itself.       %%
%%                                                                 %%
%%%%%%%%%%%%%%%%%%%%%%%%%%%%%%%%%%%%%%%%%%%%%%%%%%%%%%%%%%%%%%%%%%%%%

%%=======================================================%%
%% to print line numbers in the margin use lineno option %%
%%=======================================================%%

%% "Numbered" switches sn-basic's citation style from author-year to the
%% numbered style the manuscript text already assumes (single \cite{key}
%% calls, no \citet-style auto-generated author names).
\documentclass[lineno,pdflatex,sn-basic,Numbered]{sn-jnl}% Basic Springer Nature Reference Style, numbered citations

%%%% Standard Packages
\usepackage{graphicx}%
\usepackage{multirow}%
\usepackage{amsmath,amssymb,amsfonts}%
\usepackage[title]{appendix}%
\usepackage{booktabs}%
%%%%

%% Manuscript notes: \NOTE, \NOTEWARN, \NOTECHECK, \NOTETODO, \NOTEQUERY.
%% Add [show] to display notes in the PDF margin, or drop the option
%% entirely for a clean submission copy (both modes degrade gracefully).
\usepackage[mark]{manuscriptnotes}

\raggedbottom

\begin{document}

\title[Selection and Relevance of Radiomic Features for Reperfusion Prediction]{Selection and Relevance of Radiomic Features for Reperfusion Prediction}

%%=============================================================%%
%% GivenName	-> \fnm{Joergen W.}
%% Particle	-> \spfx{van der} -> surname prefix
%% FamilyName	-> \sur{Ploeg}
%% Suffix	-> \sfx{IV}
%%=============================================================%%

\author[1]{\fnm{Petra} \sur{Nemcekova}}
\author[1]{\fnm{Jiri} \sur{Chmelik}}
\author[1]{\fnm{Tomas} \sur{Vicar}}
\author[2]{\fnm{Mahsa} \sur{Mojtahedi}}
\author[3]{\fnm{Manon} \sur{Kappelhof}}
\author[3]{\fnm{Jan W} \sur{Hoving}}
\author[3]{\fnm{Fabiano} \sur{Cavalcante}}
\author[9,10]{\fnm{Kilian} \sur{Treurniet}}
\author[2,3]{\fnm{Nerea} \sur{Arrarte Terreros}}
\author[2,3]{\fnm{Praneeta R} \sur{Konduri}}
\author[5]{\fnm{Aad} \spfx{van der} \sur{Lugt}}
\author[6]{\fnm{Diederik} \sur{Dippel}}
\author[8]{\fnm{Jonathan} \sur{Coutinho}}
\author[11]{\fnm{Wim} \spfx{van} \sur{Zwam}}
\author[8]{\fnm{Yvo} \sur{Roos}}
\author[3]{\fnm{Bart J} \sur{Emmer}}
\author[3]{\fnm{Charles B L M} \sur{Majoie}}
\author*[1]{\fnm{Roman} \sur{Jakubicek}}\email{jakubicek@vut.cz}
\author[2,3]{\fnm{Henk A.} \sur{Marquering}}

\affil[1]{\orgdiv{Department of Biomedical Engineering and Bioinformatics}, \orgname{Brno University of Technology}, \orgaddress{\street{Technick\'a 3082/12}, \city{Brno}, \postcode{61200}, \country{Czech Republic}}}

\affil[2]{\orgdiv{Department of Biomedical Engineering and Physics}, \orgname{Amsterdam University Medical Center}, \orgaddress{\street{Meibergdreef 9}, \city{Amsterdam}, \postcode{1105 AZ}, \country{Netherlands}}}

\affil[3]{\orgdiv{Department of Radiology and Nuclear Medicine}, \orgname{Amsterdam University Medical Center}, \orgaddress{\street{Meibergdreef 9}, \city{Amsterdam}, \postcode{1105 AZ}, \country{Netherlands}}}

\affil[5]{\orgdiv{Department of Radiology and Nuclear Medicine}, \orgname{Erasmus MC}, \orgaddress{\street{Dr.\ Molewaterplein 40}, \city{Rotterdam}, \postcode{3015 GD}, \country{Netherlands}}}

\affil[6]{\orgdiv{Department of Neurology}, \orgname{Erasmus MC UMC}, \orgaddress{\street{Dr.\ Molewaterplein 40}, \city{Rotterdam}, \postcode{3015 GD}, \country{Netherlands}}}

\affil[8]{\orgdiv{Department of Neurology}, \orgname{Amsterdam UMC}, \orgaddress{\street{Meibergdreef 9}, \city{Amsterdam}, \postcode{1105 AZ}, \country{Netherlands}}}

\affil[9]{\orgdiv{Research Bureau of Radiology and Nuclear Medicine}, \orgname{Amsterdam UMC}, \orgaddress{\street{Van der Boechorststraat 6a}, \city{Amsterdam}, \postcode{1105 AZ}, \country{Netherlands}}}

\affil[10]{\orgdiv{Department of Radiology}, \orgname{The Hague Medical Center}, \orgaddress{\street{Lijnbaan 32}, \city{The Hague}, \postcode{2262 BA}, \country{Netherlands}}}

%%% Original elsarticle source labelled this author's affiliation "DRNMM",
%%% but only defined an affiliation block labelled "CARIM" -- an undefined
%%% cross-reference that would not have compiled. Mapped here to CARIM
%%% (Maastricht UMC), the one affiliation block left otherwise unused, and
%%% consistent with this author's known institution; verify before submission.
\affil[11]{\orgdiv{Department of Radiology and Nuclear Medicine, Cardiovascular Research Institute Maastricht (CARIM)}, \orgname{Maastricht UMC}, \orgaddress{\street{UNS 50, Universiteitssingel 50}, \city{Maastricht}, \postcode{6229 HX}, \country{Netherlands}}}

%%==================================%%
%% Unstructured abstract            %%
%%==================================%%

\abstract{\NOTETODO{Draft the unstructured abstract here ($\le$200 words, no citations, no equations, no subheadings). The elsarticle draft only had the four structural labels below filled with planning notes, not actual sentences -- merge them into flowing prose once the Results/Discussion sections are finalized: Background and Objectives (why eTICI$\geq$2B prediction is clinically relevant; the gap left by van~Voorst et al., 2023, who tested one selector/classifier combination); Methods (scope of the factorial design: feature-set granularity $\times$ selection method $\times$ classifier, cross-validation plus external validation, permutation test); Results (key AUC values and their significance against the permutation baseline); Conclusions (whether the choice of feature-selection method changes the clinical conclusion).}}

\keywords{brain stroke, thrombus, radiomics, CT, feature selection}

%%% The elsarticle draft also carried an Elsevier-specific "Research
%%% highlights" list (3-5 bullets, max 85 characters each), which sn-jnl has
%%% no equivalent environment for. Kept here for reference in case another
%%% Elsevier-style submission needs it later:
%%%   - Predictive Trained models
%%%   - Trained models freely available and repository
%%%   - Explore the influence of the method and dataset to prediction accuracy?
%%%   - Explore the influence of segmentation size to prediction accuracy?
%%% It also had an (empty) \begin{graphicalabstract} placeholder, likewise
%%% dropped here for the same reason.

\maketitle

\section{Tables}\label{sec1}

% Auto-generated by make_latex_tables.py - do not edit by hand.
% Individual tables also available as separate .tex files in this folder.

% ---- 4.1a thrombus vs centerline ----
\begin{table}[htbp]
\centering
\small
\caption{Whole-thrombus vs.\ centerline segmentation, tested separately within each of the 3 feature representations (voxel: raw per-voxel values; distribution: per-patient histograms; global: per-patient summary statistics), paired DeLong and McNemar tests on 141 NOIV patients (out-of-fold, pooled across 5 folds). Benjamini--Hochberg correction over $m=12$ tests; 0/12 survive.}
\label{tab:4-1a-thrombus-vs-centerline}
\begin{tabular}{lrrrrrr}
\toprule
Representation & Method & $\Delta$AUC & $p$ (DeLong) & $p$ (McNemar) & $q$ (BH, $m$=12) & Sig.\ (BH) \\
\midrule
Distribution & RF+LASSO & $-0.042$ & $0.5591$ & $\mathbf{0.0072}$ & $0.0583$ & No \\
Distribution & RF+RFI & $+0.180$ & $\mathbf{0.0094}$ & $\mathbf{0.0282}$ & $0.1125$ & No \\
Distribution & SVC+LASSO & $+0.068$ & $0.2247$ & $0.3662$ & $0.6250$ & No \\
Distribution & SVC+RFI & $+0.060$ & $0.3821$ & $0.5827$ & $0.6250$ & No \\
Global & RF+LASSO & $+0.175$ & $\mathbf{0.0047}$ & $\mathbf{0.0375}$ & $0.1125$ & No \\
Global & RF+RFI & $+0.204$ & $\mathbf{0.0002}$ & $0.2561$ & $0.6147$ & No \\
Global & SVC+LASSO & $+0.197$ & $\mathbf{0.0006}$ & $\mathbf{0.0097}$ & $0.0583$ & No \\
Global & SVC+RFI & $-0.058$ & $0.4464$ & $0.5896$ & $0.6250$ & No \\
Voxel & RF+LASSO & $+0.165$ & $\mathbf{0.0018}$ & $0.3877$ & $0.6250$ & No \\
Voxel & RF+RFI & $+0.116$ & $0.0911$ & $0.6250$ & $0.6250$ & No \\
Voxel & SVC+LASSO & $+0.068$ & $0.1575$ & $0.5488$ & $0.6250$ & No \\
Voxel & SVC+RFI & $-0.050$ & $0.5404$ & $0.5000$ & $0.6250$ & No \\
\bottomrule
\end{tabular}
\end{table}


% Po korekcii na mnohonásobné testovanie (BH, m=12) sa nenašiel štatisticky významný rozdiel medzi whole-thrombus a centerliestovaných kombinácií (3reprezentácie × 4 metódy), na hladine α=0,05.
%                                                                                   Toto je "neexistuje dôkaz rozdinie je to to isté ako "dôkaz, že rozdiel neexistuje" (evidence of absence). Test proti nulovej hypotéze "žiadny rozdiel" ti nedovolí pozitívne potvrdiť, že whole a centerline sú ekvivalentné — len ti hovorí, že dáta neposkytujú dosť silný dôkaz na zamietnutie tejto nulovej hypotézy. Keby si chcelatvrdiť "sú ekvivalentné", potret (napr. TOST, ktorý sompoužíval v sekcii 4.4), nie obyčajný test rozdielu.

% Prečo je to dôležité si všimnúť — nekorigované p-hodnoty vyzerali inak

% Toto je presne prípad, prečo korekcia zmenila záver. Bez nej by si mohla dôjsť k opačnému dojmu:

% ┌──────────────────────┬────────────┬────────┐
% │      kombinácia      │ p (DeLong) │ q (BH) │
% ├──────────────────────┼────────────┼────────┤
% │ Global/RF+RFI        │     0,0002 │  0,307 │                                    ├──────────────────────┼───────
% │ Global/SVC+LASSO     │     0,0006 │  0,058 │
% ├──────────────────────┼────────────┼────────┤
% │ Global/RF+LASSO      │     0,0047 │  0,113 │                                    ├──────────────────────┼───────
% │ Distribution/SVC+RFI │     0,0017 │  0,277 │
% └──────────────────────┴────────────┴────────┘

% Nekorigovane vyzerajú štyri z 12 kombinácií "vysoko signifikantné" (p < 0,005) — ale po zohľadnení toho, že testuješ 12 hypotéz naraz, žiadna neprežije q < 0,05. Toto je presne ten istý princíp, o ktorom sme sa bavili pri sekcii 4.2 — a tu vidíš, prečo na tom záleží: bez korekcie by si mohla mylne uzavrieť, že global-based príznaky sú citlivé na výber regiónu (whole vs centerline), keď to v skutočnosti mohol byť artefakt viacnásobného testo
% Dve kombinácie sú tesne pod hranicou                                              
% Distribution/RF+LASSO (McNemar) a Global/SVC+LASSO majú q=0,058 — veľmi blízko 0,05. Toto by som v texte článku opísal ako "trend, ktorý neprežil korekciu", nie ako nulový výsledok bez akéhokoľvek náznaku — je rozdiel medzi q=0,058 a q=0,80.

% Ešte jedna vec, ktorú stojí za to spomenúť                                        
% Všimni si Global/RF+RFI: DeLong p=0,0002 (veľmi malé), ale McNemar p=0,256 (vôbec nie malé). To nie je chyba — DeLong testuje poradie/AUC (ako dobre model rozlišuje pacientov pravdepodobnostne), McNemar tes pri prahu 0,5. Model môže maťvýrazne odlišné AUC medzi whole a centerline, a napriek tomu robiť takmer rovnaké finálne rozhodnutia (0/1) — to je vecne zaujímavé zistenie samo osebe, nie len metodická poznámka.

% Zhrnutie: tvoja formulácia je v princípe správna, len by som ju sprísnil na "nenašla som
% štatisticky významný rozdiel" niel nie je" — a do článku by som pridal poznámku o dvoch takmer-signifikantných kombináciách a o rozpore DeLong/McNemar pri Global/RF+RFI.

% ---- 4.1b Friedman ----
\begin{table}[htbp]
\centering
\small
\caption{Friedman rank test, method vs.\ method across the 6 feature sets (ranks of mean 5-fold AUC within each feature set). $F_F=0.84$ on (3, 15) df (critical value at $\alpha=0.05$: 3.29), $p=0.491$ -- omnibus not significant; no post-hoc pairwise tests performed.}
\label{tab:4-1b-friedman}
\begin{tabular}{lrr}
\toprule
Method & Mean rank & Mean AUC (across 6 feature sets) \\
\midrule
RF+LASSO & 2.50 & 0.503 \\
RF+RFI & 2.33 & 0.530 \\
SVC+LASSO & 3.17 & 0.543 \\
SVC+RFI & 2.00 & 0.488 \\
\bottomrule
\end{tabular}
\end{table}


% Inú otázku než predchádzajúce dve tabuľky — tu sa netestuje "je model lepší než náhoda", ale "záleží vôbec na tom, akú metódu (selektor+klasifikátor) si vyberieš?" Porovnáva sa RF+LASSO, RF+RFI, SVC+LASSO, SVC+RFI medzi sebou, naprieč všetkými 6 feature setmi naraz.

% Ako Friedmanov test funguje (prečo poradia, nie AUC priamo)

% Pre každý zo 6 feature setov sa 4 metódy zoradia podľa ich priemernej AUC (1 = najhoršia, 4 = najlepšia v danom feature sete), a potom sa testuje, či sú tieto poradia systematicky odlišné naprieč feature setmi, alebo len náhodná fluktuácia. Neporovnáva sa teda absolútna hodnota AUC medzi metódami (tá sa mení kombináciu od kombinácie), ale relatívne poradie — to je robustnejšie voči tomu, že rôzne feature sety majú úplne inú škálu ťažkosti.

% Interpretácia výsledku

% - Mean rank: priemernajlepšie.
% - F_F = 1,92 oproti kritickej hodnote 3,29 (α=0,05, df=(3,15)) → F_F < kritická hodnota, test nezamieta nulovú hypotézu
% - p = 0,169 → nesignifikantné
% - Žiadne post-hoc párové testy sa nerobili (správne — pri nesignifikantnom omnibuse by to bol lov na náhodné rozdiely)

% Čo to znamená

% Nenašiel sa štatisticky významný rozdiel medzi metódami — poradia, ktoré vidíš (RF+RFI najhoršie, SVC+RFI najlepšie), sú zlučiteľné s náhodnou fluktuáciou, nie s tým, že by voľba selektora/klasifikátora systematicky ovplyvňovala výkon.

% Toto je žiadaný záver (presne ako to formulovala pôvodná poznámka statistika_co_spocitat.html) — nie sklamanie. Prečo:

% - Keby na metóde výrazne záležalo, znamenalo by to, že výsledky práce sú krehké voči arbitrárnej voľbe algoritmu — čo by podkopalo dôveryhodnosť hlavného zistenia.
% - Skutočnosť, že metódy sa navzájom nelíšia, ale niektoré kombinácie feature_set×metóda predsa len prekonávajú náhodu (ako si videla v tabuľke 4.2 — 6 z 16), naznačuje, že rozdiel robí feature set (reprezentácia dát), nie voľba klasifikátora/selektora. To je vecne čistejší a silnejší príbeh pre článok, než keby výsledok závisel od toho, ktorý konkrétny algoritmus si natrénovala.

% Jedna vec na pozor pri interpretácii v texte: malá vzorka (6 "blokov" = 6 feature setov, len 4 "ošetrenia" = metódy) znamená, že tento test má nízku silu (power) — nesignifikantný výsledok pri n=6 blokov neznamená silný dôkaz ekvivalencie metód, len že sa nenašiel dostatočný dôkaz rozdielu (rovnaká výhrada ako pri predošlej tabuľke).


% ---- 4.1c CI/calibration ----

\begin{table}[htbp]
\centering
\small
\caption{Confidence intervals and calibration for the best-performing method of each feature set, pooled over 141 NOIV out-of-fold patients (DeLong CI for AUC, Wilson CI for accuracy, class-stratified bootstrap CI for balanced accuracy, $n_{\mathrm{boot}}=5000$).}
\label{tab:4-1c-ci-calibration}
\adj
\begin{tabular}{llrrrr}
\toprule
Feature set & Best method & AUC (95\% DeLong CI) & ACC (95\% Wilson CI) & Balanced ACC (95\% boot.\ CI) & Brier \\
\midrule
Distribution & RF+RFI & 0.639 (0.546--0.732) & 0.617 (0.535--0.693) & 0.583 (0.505--0.658) & 0.244 \\
Distribution (skel.) & SVC+LASSO & 0.493 (0.395--0.591) & 0.475 (0.395--0.557) & 0.473 (0.389--0.557) & 0.270 \\
Global & SVC+LASSO & 0.648 (0.557--0.739) & 0.603 (0.520--0.680) & 0.587 (0.506--0.667) & 0.266 \\
Global (skel.) & SVC+RFI & 0.574 (0.474--0.673) & 0.475 (0.395--0.557) & 0.451 (0.374--0.530) & 0.251 \\
Voxel & RF+LASSO & 0.628 (0.534--0.721) & 0.567 (0.485--0.646) & 0.490 (0.463--0.517) & 0.246 \\
Voxel (skel.) & SVC+LASSO & 0.524 (0.427--0.622) & 0.560 (0.478--0.640) & 0.489 (0.451--0.530) & 0.249 \\
\bottomrule
\end{tabular}
\end{table}

% (Drobnosť: AUC tu je 0,667 pre distribution/svc_rfi, kým v tabuľke 4.2 bolo 0,675 — nie je to chyba, len iná štatistika: tu je to jedno poolované AUC cez 141 pacientov, v 4.2 priemer 5 samostatných foldových AUC. Blízke, ale nie identické čísla sú očakávané.)

% Dôležité upozornenie — CI tu nie je hypotézový test s korekciou

% Táto tabuľka len opisuje neistotu jedného čísla (deskriptívne CI), nekoriguje na to, že si vybrala "najlepší" model z 4 kandidátov na feature_set (výberová zaujatosť) ani na 6 súbežných feature setov. Nemôžeš teda povedať "distribution je signifikantne lepší než náhoda na α=0,05" len z tejto tabuľky — na to slúži formálny test v 4.2 (ktorý navyše sedí).

% Nečakaný, ale dôležitý nález — Brierovo skóre

% Toto je vec, ktorú by som vypichol samostatne. Konštantný prediktor na skutočnej prevalencii (0,5816) dáva Brier = 0,2433 (počítal som to už v sekcii 4.3). Porovnaj:

% Ani jeden z 6 "najlepších" modelov neprekonáva Brierovým skóre triviálnu konštantu (predikovať pre každého pacienta rovnako jeho prevalenciu). To znamená: aj tam, kde model má rozumné AUC (poradie/diskriminácia je v poriadku — vie zoradiť pacientov), jeho surové pravdepodobnosti sú zle kalibrované — číselne horšie než najhlúpejší možný odhad. Toto je typický vzorec, keď model dobre rozlišuje, ale jeho výstupné skóre nie sú skutočné kalibrované pravdepodobnosti (často sa to opravuje Platt/isotonickou kalibráciou — mimochodom presne to, čo si robila pri svc_lasso v inom kontexte tejto session).

% Zhrnutie: dva nezávislé zistenia z tejto tabuľky — (1) len distribution a global (whole-thrombus) majú CI mimo náhody, čo potvrdzuje 4.2; (2) žiadny model nie je dobre kalibrovaný v absolútnom zmysle, aj tie so signálom v AUC. Do článku by druhé zistenie zaslúžilo vlastnú vetu — je to bežný, ale dôležitý caveat pre klinickú interpretáciu pravdepodobností.


% ---- 4.2 permutation test ----
\begin{table}[htbp]
\centering
\small
\caption{Permutation test against the label-shuffled null (5000 shared permutations per fold; test statistic = mean AUC across the 5 folds), $m=16$ aggregate feature-set combinations (voxel excluded). $\Delta\mathrm{AUC}$ and $\kappa$ relative to the null's own mean, not to 0.5. 4/16 combinations survive Benjamini--Hochberg correction.}
\label{tab:4-2-permutation}
\begin{tabular}{lrrrrrrr}
\toprule
Feature set & Method & AUC & $\Delta$AUC & $\kappa$ & $p$ & $q$ (BH, $m$=16) & Sig.\ (BH) \\
\midrule
Global & SVC+LASSO & 0.658 & $+0.159$ & 0.317 & $\mathbf{0.0012}$ & $\mathbf{0.0107}$ & Yes \\
Global & RF+LASSO & 0.646 & $+0.147$ & 0.294 & $\mathbf{0.0016}$ & $\mathbf{0.0107}$ & Yes \\
Distribution & RF+RFI & 0.645 & $+0.146$ & 0.291 & $\mathbf{0.0020}$ & $\mathbf{0.0107}$ & Yes \\
Global & RF+RFI & 0.617 & $+0.117$ & 0.233 & $\mathbf{0.0108}$ & $\mathbf{0.0432}$ & Yes \\
Global (skel.) & SVC+RFI & 0.568 & $+0.069$ & 0.137 & $0.0850$ & $0.2698$ & No \\
Distribution & SVC+LASSO & 0.560 & $+0.060$ & 0.120 & $0.1012$ & $0.2698$ & No \\
Global & SVC+RFI & 0.516 & $+0.016$ & 0.031 & $0.3881$ & $0.8871$ & No \\
Distribution & SVC+RFI & 0.487 & $-0.012$ & -0.025 & $0.5897$ & $0.9972$ & No \\
Distribution (skel.) & SVC+LASSO & 0.482 & $-0.019$ & -0.037 & $0.6433$ & $0.9972$ & No \\
Global (skel.) & RF+LASSO & 0.466 & $-0.034$ & -0.067 & $0.7429$ & $0.9972$ & No \\
Distribution (skel.) & RF+RFI & 0.461 & $-0.038$ & -0.075 & $0.7658$ & $0.9972$ & No \\
Global (skel.) & SVC+LASSO & 0.453 & $-0.047$ & -0.094 & $0.8178$ & $0.9972$ & No \\
Distribution (skel.) & RF+LASSO & 0.468 & $-0.032$ & -0.064 & $0.8836$ & $0.9972$ & No \\
Distribution (skel.) & SVC+RFI & 0.429 & $-0.072$ & -0.143 & $0.9152$ & $0.9972$ & No \\
Global (skel.) & RF+RFI & 0.396 & $-0.105$ & -0.210 & $0.9838$ & $0.9972$ & No \\
Distribution & RF+LASSO & 0.353 & $-0.147$ & -0.293 & $0.9972$ & $0.9972$ & No \\
\bottomrule
\end{tabular}
\end{table}


% Toto je hlavný pozitívny test celej štúdie: pre každú z 16 kombinácií (4 whole-thrombus feature sety × 4 metódy — skeleton varianty zámerne vylúčené z hlavnej rodiny, pozri predchádzajúcu diskusiu) sa testuje, či reálne natrénovaný model predikuje lepšie, než keby sa pacientske labely náhodne prehádzali.

% Ako sa to počíta

% 1. Pre každú kombináciu, pre každý z 5 foldov: 5000× sa náhodne prehádžu labely, model sa nepretrénováva (len sa prehodnotí na tých istých predikciách), spočíta sa AUC — vznikne nulové rozdelenie "čo by dosiahla náhoda".
% 2. Testová štatistika = priemer AUC cez 5 foldov (reálny aj nulový, párovaný podľa čísla iterácie).
% 3. p-hodnota = podiel z 5000 permutácií, ktoré dosiahli aspoň takú vysokú priemernú AUC ako reálny model.
% 4. q = p po BH korekcii cez m=16.

% Ako čítať stĺpce

% - AUC — priemerná AUC reálneho modelu cez 5 foldov
% - ΔAUC — o koľko je to nad priemerom nulového rozdelenia tejto konkrétnej kombinácie (nie nad 0,5 — nula sa mierne líši kombináciu od kombinácie)
% - κ — normalizovaná verzia ΔAUC, ΔAUC / (1 − null_mean) — rovnaký zápis funguje aj pre ACC, kde by odčítanie od 0,5 bolo zavádzajúce
% - q (BH, m=16) — rozhodujúci stĺpec

% Výsledok

% 6 z 16 kombinácií prežíva korekciu (q < 0,05) — jediné miesto v celej analýze 4.1–4.4, kde niečo formálne prežije:

% ┌──────────────┬───────────┬───────┬───────┐
% │ feature set  │  metóda   │  AUC  │   q   │
% ├──────────────┼───────────┼───────┼───────┤
% │ Distribution │ SVC+RFI   │ 0,675 │ 0,008 │
% ├──────────────┼───────────┼───────┼───────┤
% │ Global       │ SVC+LASSO │ 0,658 │ 0,008 │
% ├──────────────┼───────────┼───────┼───────┤
% │ Global       │ RF+LASSO  │ 0,646 │ 0,008 │
% ├──────────────┼───────────┼───────┼───────┤
% │ Distribution │ RF+RFI    │ 0,645 │ 0,008 │
% ├──────────────┼───────────┼───────┼───────┤
% │ Global       │ SVC+RFI   │ 0,627 │ 0,024 │
% ├──────────────┼───────────┼───────┼───────┤
% │ Global       │ RF+RFI    │ 0,617 │ 0,029 │
% └──────────────┴───────────┴───────┴───────┘

% Zvyšných 10 kombinácií — všetky skeleton varianty plus Distribution/SVC+LASSO (q=0,231) — nemá dôkaz signálu; AUC sa pohybuje okolo alebo pod 0,5.

% Tri veci, ktoré z tabuľky vyplývajú

% 1. Signál je len vo whole-thrombus dátach, nie v skeleton. Všetkých 6 signifikantných kombinácií pochádza z ne-skeletonizovaných feature setov. Toto priamo súvisí s tabuľkou 4.1a, kde whole-vs-centerline rozdiel formálne neprežil korekciu, ale dve kombinácie (Distribution/RF+LASSO, Global/SVC+LASSO) boli tesne pod hranicou q=0,058 — konzistentné s tým, že centerline dáta majú slabší signál.

% 2. Metóda (selektor/klasifikátor) nerozhoduje o tom, či je signál — rozhoduje o tom feature set. To sedí s tabuľkou 4.1b (Friedman): žiadny rozdiel medzi RF/SVC, LASSO/RFI naprieč feature setmi. Tu vidíš prečo — v rámci global aj distribution sú signifikantné 3 zo 4 metód (chýba len SVC+LASSO pri distribution, RF+RFI je najhraničnejšia pri global). Signál teda nesie feature set, nie konkrétny algoritmus.

% 3. Distribution/RF+LASSO je pozoruhodne zlý — AUC 0,353, ΔAUC = −0,147 (hlboko pod vlastnou náhodnou nulou). Model systematicky predikuje opačným smerom, než by mal. Test je jednostranný (testuje len "lepšie než náhoda"), takže sa to formálne neoznačí ako "signifikantné", ale je to vecne zaujímavý nález — táto konkrétna kombinácia je nestabilná/nespoľahlivá, nie len "bez signálu".

% Zhrnutie pre článok: existuje reálny, korigovaný pozitívny nález — ale obmedzený na whole-thrombus distribution a global feature sety (6/16 kombinácií), nie na skeleton varianty ani konkrétnu voľbu klasifikátora.

% ---- 4.3a ours vs reference ----
\begin{table}[htbp]
\centering
\small
\caption{Our pipeline vs.\ the van~Voorst/Henks exact-replication reference pipeline, paired DeLong test on the same 141 real outer-test-fold patients, $m=16$. 0/16 combinations survive Benjamini--Hochberg correction. The headline accuracy advantage (up to $+0.11$) collapses to at most $+0.035$ once the single driving fold (split index 1) is excluded (not shown in this table; see \texttt{stats\_4\_3\_ours\_vs\_reference.csv}).}
\label{tab:4-3a-ours-vs-reference}
\begin{tabular}{lrrrrrrr}
\toprule
Feature set & Method & AUC (ours) & AUC (ref.) & $\Delta$AUC & $p$ & $q$ (BH, $m$=16) & Sig.\ (BH) \\
\midrule
Global & SVC+LASSO & 0.648 & 0.492 & $+0.157$ & $\mathbf{0.0224}$ & $0.1519$ & No \\
Global & RF+LASSO & 0.641 & 0.492 & $+0.149$ & $\mathbf{0.0329}$ & $0.1519$ & No \\
Distribution & RF+LASSO & 0.360 & 0.492 & $-0.131$ & $\mathbf{0.0357}$ & $0.1519$ & No \\
Distribution & RF+RFI & 0.639 & 0.492 & $+0.147$ & $\mathbf{0.0380}$ & $0.1519$ & No \\
Global & RF+RFI & 0.606 & 0.492 & $+0.114$ & $0.1002$ & $0.3207$ & No \\
Distribution (skel.) & RF+LASSO & 0.402 & 0.492 & $-0.090$ & $0.1574$ & $0.4198$ & No \\
Global (skel.) & RF+RFI & 0.402 & 0.492 & $-0.090$ & $0.1921$ & $0.4391$ & No \\
Global (skel.) & SVC+RFI & 0.574 & 0.492 & $+0.082$ & $0.2511$ & $0.5022$ & No \\
Distribution & SVC+LASSO & 0.562 & 0.492 & $+0.070$ & $0.3453$ & $0.6081$ & No \\
Distribution (skel.) & SVC+RFI & 0.429 & 0.492 & $-0.063$ & $0.3800$ & $0.6081$ & No \\
Global (skel.) & SVC+LASSO & 0.451 & 0.492 & $-0.041$ & $0.5543$ & $0.8062$ & No \\
Distribution (skel.) & RF+RFI & 0.459 & 0.492 & $-0.033$ & $0.6170$ & $0.8226$ & No \\
Global (skel.) & RF+LASSO & 0.465 & 0.492 & $-0.026$ & $0.7060$ & $0.8308$ & No \\
Global & SVC+RFI & 0.516 & 0.492 & $+0.024$ & $0.7270$ & $0.8308$ & No \\
Distribution & SVC+RFI & 0.489 & 0.492 & $-0.003$ & $0.9694$ & $0.9837$ & No \\
Distribution (skel.) & SVC+LASSO & 0.493 & 0.492 & $+0.001$ & $0.9837$ & $0.9837$ & No \\
\bottomrule
\end{tabular}
\end{table}


% Iná otázka než 4.2 — tam sa porovnávalo "náš model vs náhoda", tu sa porovnáva "náš model vs van Voorstova/Henkova referenčná pipeline" (jeden pevný ExtraTrees model, žiadny výber príznakov, natrénovaný na fixnej sade "rv"). Test je párovaný DeLong na tých istých 141 pacientoch zo skutočných testovacích foldov (nie z vnútorných holdoutov, čo bola pôvodná chyba obrázka, ktorú sme opravovali skôr).

% Prečo je "AUC (ref.)" v každom riadku rovnaké číslo (0,492)

% Toto nie je chyba — je to tá istá jedna referenčná pipeline, porovnávaná postupne proti každej z 16 našich kombinácií. Referencia sama osebe má jednu pevnú (poolovanú) AUC = 0,492, nezávislú od toho, s ktorým naším modelom sa práve porovnáva.

% Výsledok

% 0 z 16 prežíva BH korekciu (min q = 0,122) — presne to, čo si videla aj v súhrnnej tabuľke. Ale všimni si, že 5 kombinácií malo surové p < 0,05:

% ┌──────────────┬───────────┬────────┬───────┬───────┐
% │ feature set  │  metóda   │  ΔAUC  │   p   │   q   │
% ├──────────────┼───────────┼────────┼───────┼───────┤
% │ Distribution │ SVC+RFI   │ +0,175 │ 0,012 │ 0,122 │
% ├──────────────┼───────────┼────────┼───────┼───────┤
% │ Global       │ SVC+LASSO │ +0,157 │ 0,022 │ 0,122 │
% ├──────────────┼───────────┼────────┼───────┼───────┤
% │ Global       │ RF+LASSO  │ +0,149 │ 0,033 │ 0,122 │
% ├──────────────┼───────────┼────────┼───────┼───────┤
% │ Distribution │ RF+LASSO  │ −0,131 │ 0,036 │ 0,122 │
% ├──────────────┼───────────┼────────┼───────┼───────┤
% │ Distribution │ RF+RFI    │ +0,147 │ 0,038 │ 0,122 │
% └──────────────┴───────────┴────────┴───────┴───────┘

% Bez korekcie by si mohla povedať "5 z 16 modelov je lepších než referencia" — po korekcii na 16 súbežných testov nič neprežije.

% Prečo to nesedí s tabuľkou 4.2 (kde 6/16 prežilo)

% Toto je dôležitý bod. Rovnaké najlepšie kombinácie (Distribution/SVC+RFI, Global/SVC+LASSO, Global/RF+LASSO, Distribution/RF+RFI) tu majú horšie p-hodnoty než v 4.2, hoci porovnávaš prakticky to isté (dobrý model vs. niečo blízke náhode — referencia má AUC 0,492, teda skoro presne 0,5). Dôvod:

% - V 4.2 sa porovnávalo proti permutačnej nule s 5000 vzorkami — veľmi presne odhadnutá, málo šumu.
% - Tu sa porovnáva proti jednému konkrétnemu modelu s vlastnou variabilitou medzi foldmi — DeLongov test musí zohľadniť neistotu oboch kriviek naraz (kovariančný člen), a referenčná pipeline má medzi foldmi značne nerovnomerný výkon (spomeň si, jeden fold mal AUC 0,54, ostatné ~0,48–0,52). Táto extra neistota z referencie robí test menej silný (nižšia štatistická sila), aj keď smer rozdielu je rovnaký.

% Praktický záver: "náš model je lepší než náhoda" je štatisticky podložené tvrdenie (4.2). "Náš model je lepší než referenčná pipeline" nie je — napriek tomu, že čísla vyzerajú priaznivo (ΔAUC až +0,175).

% Distribution/RF+LASSO — opäť problematická kombinácia

% Rovnaká kombinácia, ktorá v 4.2 mala AUC hlboko pod náhodou (0,353), tu má záporné Δ voči referencii (−0,131, p=0,036) — teda je horšia aj než skoro-náhodná referenčná pipeline. Toto je tretí nezávislý test (po 4.2), kde sa táto kombinácia javí ako systematicky nespoľahlivá — stojí za samostatnú vetu v diskusii, nie len za tichý riadok v tabuľke.

% Poznámka z popisu tabuľky — jeden fold ťahá celý rozdiel

% Caption pripomína: pôvodná výhoda v ACC (až +0,11) klesne na max. +0,035, keď sa vylúči jeden konkrétny fold (split index 1). To znamená, že aj to skromné, nekorigované zdanie výhody nad referenciou z veľkej časti stojí na jednom foldu — ďalší dôvod, prečo záver "sme lepší než referencia" v tejto podobe dát neobstojí.

% Zhrnutie pre článok: napriek priaznivo vyzerajúcim rozdielom AUC/ACC, žiadny formálny, korigovaný dôkaz, že naša pipeline prekonáva referenčnú van Voorstovu pipeline, neexistuje — a zdanlivá výhoda navyše závisí najmä od jedného foldu.

% ---- 4.3b reference calibration ----
\begin{table}[htbp]
\centering
\small
\caption{Calibration of the reference pipeline on 141 NOIV out-of-fold patients. The reference is worse than a constant predictor at the observed prevalence; its predicted probabilities vary too little (SD $=0.022$) for the calibration slope/intercept to be meaningfully interpreted.}
\label{tab:4-3b-reference-calibration}
\begin{tabular}{lr}
\toprule
Quantity & Value \\
\midrule
Brier score (reference) & 0.2468 \\
Brier score, constant $\hat p=0.5$ & 0.2500 \\
Brier score, constant $\hat p=$ prevalence & 0.2433 \\
Predicted-probability range & [0.4983, 0.5570] \\
Predicted-probability SD & 0.0219 \\
Calibration slope / intercept & 2.16 / 0.24 \\
\bottomrule
\end{tabular}
\end{table}


% Toto je vysvetlenie, prečo referenčná pipeline v tabuľke 4.3a vyšla s AUC len 0,492 (prakticky náhoda). Kým 4.3a testovala diskrimináciu (vie model zoradiť pacientov podľa rizika?), táto tabuľka testuje kalibráciu (sú jeho pravdepodobnosti vôbec zmysluplné čísla?) — to sú dve nezávislé vlastnosti modelu, ktoré sa dajú pokaziť samostatne.

% Kľúčové zistenie — Brierovo skóre

% Brierovo skóre = priemer (predikovaná pravdepodobnosť − skutočný výsledok)² cez 141 pacientov. Nižšie = lepšie, 0 = dokonalé.

% ┌──────────────────────────────────┬────────┐
% │                                  │ Brier  │
% ├──────────────────────────────────┼────────┤
% │ Referenčná pipeline              │ 0,2468 │
% ├──────────────────────────────────┼────────┤
% │ Konštanta 0,5 (hodenie mincou)   │ 0,2500 │
% ├──────────────────────────────────┼────────┤
% │ Konštanta na prevalencii (0,582) │ 0,2433 │
% │ Konštanta 0,5 (hodenie mincou)   │ 0,2500 │
% ├──────────────────────────────────┼────────┤
% │ Konštanta na prevalencii (0,582) │ 0,2433 │
% └──────────────────────────────────┴────────┘

% Referenčná pipeline je horšia než najhlúpejší možný odhad — teda horšia, než keby si pre každého pacienta bez ohľadu na jeho dáta predikovala rovnaké číslo (prevalenciu ochorenia v populácii). Je len minimálne (0,0032) lepšia než hodenie mincou. Toto nie je len "slabý model" — je to model, ktorý neprináša žiadnu použiteľnú informáciu navyše oproti tomu, čo už vieš bez toho, aby si sa vôbec pozrela na dáta pacienta.

% Prečo — rozsah predikovaných pravdepodobností

% [0,4983, 0,5570], smerodajná odchýlka 0,022. Model prakticky vždy hovorí "približne 50 na 50", bez ohľadu na konkrétneho pacienta — celý rozsah predikcií naprieč všetkými 141 pacientmi je len 0,057 široký. Toto priamo vysvetľuje AUC≈0,5 z tabuľky 4.3a: keď model dáva takmer identické skóre každému, nemá z čoho zoradiť pacientov.

% Kalibračný sklon/intercept — prečo ich neinterpretovať

% Hodnoty 2,16 / 0,24 by za normálnych okolností znamenali: ideálna kalibrácia = sklon 1, intercept 0; sklon > 1 naznačuje, že model je "príliš konzervatívny" (jeho pravdepodobnosti sú stlačené bližšie k 0,5, než by mali byť). Ale tu je tento výpočet metodicky pochybný — kalibračná regresia sa fituje na vstupe, ktorý má smerodajnú odchýlku len 0,022 (takmer konštantný). Fitovanie sklonu na takmer degenerovanom prediktore je numericky nestabilné a číslo 2,16 nemá spoľahlivú interpretačnú hodnotu — je to artefakt pokusu preložiť priamku takmer jedným bodom, nie skutočný kalibračný nález.

% Zhrnutie pre článok

% Toto potvrdzuje presne to, čo navrhovala pôvodná poznámka (statistika_co_spocitat.html): správna formulácia nie je "referenčná pipeline bola horšia", ale "referenčná pipeline nekonvergovala k ničomu netriviálnemu" — jej výstup je prakticky konštantný, nie zle kalibrovaný v bežnom zmysle. Toto je dôležité rozlíšenie: chráni to pred obvinením z nefér reprodukcie (že sme referenčnú metódu implementovali zle) — problém nie je v tom, ako sme ju spustili, ale v tom, že sa sama naučila degenerované riešenie.


% ---- 4.4 Registry vs chance ----
\begin{table}[htbp]
\centering
\small
\caption{External validation on the MR~CLEAN Registry common cohort ($n=1635$, every model on identical patients). $^{\dagger}$ pre-specified primary model (highest mean 5-fold NOIV CV AUC), fixed before looking at the Registry. Analytic chance null $\mathrm{Var}_0=(n+1)/(12\,n_+n_-)$; TOST equivalence margin $0.05$ AUC. 0/16 survive Benjamini--Hochberg ($m=16$); 14/16 are TOST-equivalent to chance. Majority-class baseline accuracy is 0.553; no combination exceeds it.}
\label{tab:4-4-registry}
\begin{tabular}{lrrrrrr}
\toprule
Feature set / method & AUC & 95\% DeLong CI & $p$ (vs.\ chance) & $q$ (BH, $m$=16) & TOST-equiv. & ACC \\
\midrule
Distribution / SVC+RFI & 0.531 & 0.503--0.560 & $\mathbf{0.0283}$ & $0.4290$ & No & 0.553 \\
Global / RF+RFI & 0.472 & 0.444--0.500 & $0.0536$ & $0.4290$ & No & 0.491 \\
\textbf{Global / SVC+LASSO}$^{\dagger}$ & 0.479 & 0.451--0.508 & $0.1521$ & $0.7726$ & Yes & 0.492 \\
Distribution / RF+RFI & 0.482 & 0.453--0.510 & $0.1986$ & $0.7726$ & Yes & 0.448 \\
Global / RF+LASSO & 0.484 & 0.456--0.512 & $0.2719$ & $0.7726$ & Yes & 0.486 \\
Distribution (skel.) / RF+LASSO & 0.514 & 0.493--0.535 & $0.3456$ & $0.7726$ & Yes & 0.447 \\
Global (skel.) / RF+LASSO & 0.487 & 0.470--0.504 & $0.3785$ & $0.7726$ & Yes & 0.532 \\
Distribution (skel.) / SVC+RFI & 0.489 & 0.462--0.515 & $0.4413$ & $0.7726$ & Yes & 0.526 \\
Global (skel.) / SVC+RFI & 0.510 & 0.482--0.538 & $0.4912$ & $0.7726$ & Yes & 0.492 \\
Distribution (skel.) / RF+RFI & 0.491 & 0.463--0.519 & $0.5182$ & $0.7726$ & Yes & 0.443 \\
Global (skel.) / RF+RFI & 0.508 & 0.480--0.537 & $0.5552$ & $0.7726$ & Yes & 0.467 \\
Distribution (skel.) / SVC+LASSO & 0.507 & 0.479--0.535 & $0.6055$ & $0.7726$ & Yes & 0.544 \\
Distribution / SVC+LASSO & 0.494 & 0.466--0.522 & $0.6626$ & $0.7726$ & Yes & 0.469 \\
Distribution / RF+LASSO & 0.494 & 0.480--0.509 & $0.7010$ & $0.7726$ & Yes & 0.447 \\
Global (skel.) / SVC+LASSO & 0.505 & 0.477--0.533 & $0.7243$ & $0.7726$ & Yes & 0.505 \\
Global / SVC+RFI & 0.500 & 0.500--0.500 & $1.0000$ & $1.0000$ & Yes & 0.553 \\
\bottomrule
\end{tabular}
\end{table}


% Toto je najtvrdšia skúška v celej práci (presne tak to formulovala aj pôvodná poznámka): generalizuje sa čokoľvek natrénované na NOIV kohorte (141 pacientov) na skutočne externú, nezávislú kohortu MR CLEAN Registry? n = 1635 pacientov — všetkých 16 modelov na identických pacientoch (to je tá "common cohort" práca, ktorú sme robili na začiatku tejto session, presne kvôli tomu, aby toto porovnanie bolo férové).

% Kľúčový bod — predem určený model

% Riadok s † (Distribution/SVC+RFI) je vopred deklarovaný primárny model — vybraný pred pozretím na Registry výsledky, na základe najvyššej priemernej CV AUC na NOIV (0,675, presne ten istý model, čo bol na vrchole tabuľky 4.2). Toto je metodicky dôležité: keby si vyberala "najlepší" model až po tom, čo vidíš Registry výsledky, bolo by to data snooping — falošne optimistický výber.

% Jeho výkon na Registry: AUC = 0,481 (CI 0,453–0,510) — prakticky presne náhoda, dokonca mierne pod 0,5. Model, ktorý na NOIV vyzeral najlepšie (0,675), na skutočne nezávislých dátach nefunguje vôbec.

% Formálny test proti náhode

% p = 0,195 pre primárny model — nesignifikantné aj bez korekcie. Po BH korekcii cez všetkých 16: 0/16 prežíva, min q = 0,532 — najslabší (najviac "nulový") výsledok zo všetkých štyroch sekcií 4.1–4.4.

% TOST — pozitívny dôkaz, že je to naozaj náhoda (nie len chýbajúci dôkaz opaku)

% Toto je rozdiel oproti predchádzajúcim tabuľkám, ktoré vedeli povedať len "nenašli sme rozdiel" (absence of evidence). TOST test tu robí niečo silnejšie: aktívne potvrdzuje, že skutočná AUC leží v pásme ±0,05 okolo 0,5.

% 15 z 16 kombinácií je TOST-ekvivalentných s náhodou — teda pre 15 z 16 modelov máš pozitívny dôkaz, že ich výkon na Registry je prakticky na úrovni náhody, nie len "nedokázali sme, že nie je".

% Jedna kombinácia (Global/RF+RFI, prvý riadok) neprešla ani jedným testom — ani proti náhode (p=0,054, tesne nad hranicou), ani TOST ekvivalenciou. Táto je štatisticky neurčená — dáta nedovoľujú ani potvrdiť, že má signál, ani potvrdiť, že je ekvivalentná náhode. Pozor na interpretáciu: jej AUC = 0,472, teda pod 0,5 — to p=0,054 nie je "skoro dôkaz reálneho signálu", je to skoro dôkaz, že model predikuje systematicky opačne (rovnaký vzorec, aký sme videli pri Distribution/RF+LASSO v 4.2 a 4.3a).

% Presvedčivý, jednoduchý fakt — porovnanie s väčšinovou triedou

% Väčšinová trieda v Registry má podiel 0,553. Ani jedna z 16 kombinácií túto hodnotu nepresahuje (najlepšia ACC = 0,544). To znamená: triviálny klasifikátor, ktorý by pre každého pacienta predikoval tú istú (väčšinovú) triedu bez toho, aby sa vôbec pozrel na dáta, by porazil všetkých 16 natrénovaných modelov na presnosti.

% Zhrnutie pre článok

% Toto je rozhodujúci negatívny výsledok: model vybraný vopred podľa najlepšej CV výkonnosti na vývojovej kohorte zlyháva na skutočne nezávislej externej kohorte — nielen že neprežije korekciu, ale TOST pozitívne potvrdzuje, že jeho výkon je na úrovni náhody, a ani jeden model neprekoná ani triviálny "vždy predikuj väčšinu" prístup. V kombinácii s tabuľkou 4.3b (referenčná pipeline tiež zlyháva, len iným spôsobom — degenerovanou kalibráciou) to vytvára konzistentný obraz: čokoľvek natrénované na tejto malej NOIV kohorte (vlastná pipeline aj referenčná) negeneralizuje na externú populáciu.


% ---- Overall summary ----
\begin{table}[htbp]
\centering
\small
\caption{Summary of all statistical analyses across manuscript Results sections 4.1--4.4. $q$ = Benjamini--Hochberg-adjusted $p$-value (minimum across the family where applicable). The only finding that survives multiplicity correction is the permutation test in Section~4.2 (6/16 combinations); every comparison against the reference pipeline (4.3) and every external validation result on the Registry (4.4) is consistent with chance after correction.}
\label{tab:results-summary}
\begin{tabular}{llllllp{4.2cm}}
\toprule
Sec. & Analysis & Test & Scope & Key statistic & $p$/$q$ & Verdict \\
\midrule
4.1 & Whole vs.\ centerline (3 representations) & DeLong / McNemar & $m=12$ & 0/12 (McNemar, BH) & 0.058 & No effect survives correction \\
4.1 & Method vs.\ method (selector/classifier) & Friedman omnibus & $k=4$, $n=6$ & $F_F=0.84$ & 0.491 & Not significant (as expected) \\
4.1 & Best-model CI / calibration & DeLong / Wilson / Brier CI & $n=141$ & AUC 0.648 & -- & Best: Global/SVC+LASSO, CI [0.56, 0.74] \\
4.2 & Performance vs.\ permutation null & Permutation test (5000 perm.) & $m=16$ & 4/16 survive BH & 0.0107 & \textbf{4/16 significant} -- main positive finding \\
4.3 & Our pipeline vs.\ reference pipeline & Paired DeLong & $m=16$, $n=141$ & 0/16 survive BH & 0.152 & No difference survives correction; headline gap driven by 1 fold \\
4.3 & Reference pipeline calibration & Brier score & $n=141$ & Brier 0.247 & -- & Worse than constant-at-prevalence predictor (0.243) \\
4.4 & Pre-specified model on Registry & Analytic AUC null / TOST & $n=1635$ & AUC 0.479 & 0.152 & TOST-equivalent to chance ($p=0.020$) \\
4.4 & All 16 combinations on Registry & Analytic AUC null, BH & $m=16$, $n=1635$ & 0/16 survive BH & 0.429 & 14/16 TOST-equivalent to chance; none beat majority class \\
\bottomrule
\end{tabular}
\end{table}







\begin{appendices}





\end{appendices}


\end{document}
